\documentclass[10pt]{beamer}
\usepackage[utf8]{inputenc}
\usepackage{xeCJK}
\usepackage{graphicx}
\usepackage{mathtools}
\usepackage{amsbsy, amsthm, amsmath, amssymb, amsfonts, bbm, mathrsfs, mathtools} % math
% Allows the use of \toprule, \midrule and \bottomrule in tables
\usepackage{booktabs} 
\usepackage[most]{tcolorbox}  
\usepackage{utopia} %font utopia imported
\usetheme{Berlin}
\usecolortheme{dolphin}
% set colors
\definecolor{myNewColorA}{RGB}{126,12,110}
\definecolor{myNewColorB}{RGB}{165,85,154}
\definecolor{myNewColorC}{RGB}{203,158,197}
\definecolor{myNewColorD}{RGB}{255,192,192}
\definecolor{myNewColorE}{RGB}{29, 113, 81}
\definecolor{myNewColorF}{RGB}{245, 245, 239}
\setbeamercolor*{palette primary}{bg=myNewColorE}
\setbeamercolor*{palette secondary}{bg=myNewColorD, fg = white}
\setbeamercolor*{palette tertiary}{bg=myNewColorE, fg = myNewColorD}
\setbeamercolor*{titlelike}{fg=myNewColorE}
\setbeamercolor*{title}{bg=myNewColorE, fg = myNewColorD}
\setbeamercolor*{item}{fg=myNewColorD}
\setbeamercolor*{caption name}{fg=myNewColorE}
\usefonttheme{professionalfonts}
\usepackage{natbib}
\usepackage{hyperref}
%------------------------------------------------------------
%\titlegraphic{%\includegraphics[height=1.5cm]{nku-logo.eps}}
\setbeamerfont{title}{size=\large}
\setbeamerfont{subtitle}{size=\small}
\setbeamerfont{author}{size=\small}
\setbeamerfont{date}{size=\small}
\setbeamerfont{institute}{size=\small}
\title[University of La Habana]{Un Generador de Problemas Prueba para Optimización Binivel}
%\subtitle{The Top Forwards in the World}
\author[Francisco Vicente Suárez Bellón]{Francisco Vicente Suárez Bellón\\[10mm]{\small Tutora: Dra. C. Gemayqzel Bouza Allende }}
\date{Febrero, 2025}
%------------------------------------------------------------

\begin{document}

% The next statement creates the title page.
\frame{\titlepage}

\begin{frame}
\frametitle{Contenido}
\tableofcontents
\end{frame}

%----------------------------------------------------------------------------------------
%	INTRODUCCIÓN
%----------------------------------------------------------------------------------------
\section{Introducción}
\begin{frame}
\frametitle{Introducción}
La optimización binivel es un área fundamental en investigación operativa y teoría de juegos. Este trabajo propone:
\begin{itemize}
    \item Un generador de problemas binivel cuyos puntos estacionarios tienen propiedades específicas.
    \item Evaluar algoritmos MPEC (Mathematical Programs with Equilibrium Constraints).
\end{itemize}
El problema general se define como:
\begin{equation*}
\min_{x} F(x, y) \quad \text{s.a.} \quad x \in T, \quad y \in S(x),
\end{equation*}
donde $S(x)$ es el conjunto de soluciones óptimas del problema del nivel inferior.
\end{frame}

%----------------------------------------------------------------------------------------
%	MARCO TEÓRICO
%----------------------------------------------------------------------------------------
\section{Marco Teórico}
\begin{frame}
\frametitle{Transformación KKT}
Reemplazando el problema del nivel inferior por sus condiciones KKT:
\begin{align*}
\min_{x, y, \lambda_j} & \quad F(x, y) \\
\text{s.a.} & \quad g_i(x, y) \leq 0, \quad i = 1, \dots, q, \\
& \quad \nabla_y f(x, y) + \sum_{j=1}^s \nabla_y v_j(x, y) \lambda_j = 0, \\
& \quad v_j(x, y) \leq 0, \quad j = 1, \dots, s, \\
& \quad \lambda_j \geq 0, \quad j = 1, \dots, s, \\
& \quad v_j(x, y) \lambda_j = 0, \quad j = 1, \dots, s.
\end{align*}
\end{frame}

%----------------------------------------------------------------------------------------
%	METODOLOGÍA
%----------------------------------------------------------------------------------------
\section{Metodología}
\begin{frame}
\frametitle{Fases del Proceso}
\begin{enumerate}
    \item Obtención de puntos mínimos iniciales.
    \item Modificación de problemas originales para garantizar estacionariedad.
    \item Evaluación de algoritmos en problemas modificados.
\end{enumerate}
Métodos de reformulación:
\begin{itemize}
    \item \textbf{Big-M}: Introduce variables binarias.
    \item \textbf{SOS1}: Utiliza conjuntos especiales ordenados.
    \item \textbf{ProductMode}: Aproxima restricciones de complementariedad.
\end{itemize}
\end{frame}

%----------------------------------------------------------------------------------------
%	RESULTADOS EXPERIMENTALES
%----------------------------------------------------------------------------------------
\section{Resultados Experimentales}
\begin{frame}
\frametitle{Experimentación}
\begin{itemize}
    \item Problemas lineales: Mejora en algunos casos con Big-M.
    \item Problemas cuadráticos: Comportamiento similar entre métodos.
    \item Problemas no convexos: Coincidencia de variables del nivel superior en puntos óptimos.
\end{itemize}
\end{frame}

%----------------------------------------------------------------------------------------
%	CONCLUSIONES
%----------------------------------------------------------------------------------------
\section{Conclusiones y Trabajo Futuro}
\begin{frame}
\frametitle{Conclusiones}
\begin{itemize}
    \item El generador permite evaluar el desempeño de algoritmos en problemas específicos.
    \item Los problemas no convexos presentan mayores desafíos estructurales.
\end{itemize}
Trabajo futuro:
\begin{itemize}
    \item Ampliar la experimentación a más categorías de problemas.
    \item Investigar criterios de parada basados en puntos estacionarios.
    \item Desarrollar una interfaz gráfica más intuitiva.
\end{itemize}
\end{frame}

\end{document}